%% =============================================================================
%% chapters/06-deployment-ha.tex
%% ArchiMate Physical & High Availability Deployment Architecture Layer
%% =============================================================================

\chapter{Physical \& High Availability Deployment}
\label{chap:deployment_ha}

The Physical and Deployment Architecture layer specifies the high-availability infrastructure topology, Apache ActiveMQ Artemis multi-broker replication groups, Infinispan cache grid clustering, Docker Compose deployment configurations, and failover/recovery mechanics.

\begin{figure}[htbp]
    \centering
    \begin{adjustbox}{max width=\textwidth}
        %% =============================================================================
%% diagrams/fig-deployment-ha.tex
%% ArchiMate Physical Deployment & Artemis HA Replication Topology
%% =============================================================================
\begin{tikzpicture}[node distance=1.0cm and 0.8cm, auto]

    % --- Background Plates ---
    \draw[archi-plate-bg] (-0.8, 6.2) rectangle (15.8, 8.8);
    \node[archi-plate-title] at (-0.6, 8.6) {Client Connection Pool \& HA Topology Resolver};

    \draw[archi-plate-solid] (-0.8, 0.4) rectangle (7.2, 5.8);
    \node[archi-plate-title] at (-0.6, 5.6) {Replication Group A (group-a)};

    \draw[archi-plate-solid] (7.8, 0.4) rectangle (15.8, 5.8);
    \node[archi-plate-title] at (8.0, 5.6) {Replication Group B (group-b)};

    % Client Tier
    \node[archi-deployment-node, minimum width=14.5cm] (client_pool) at (7.5, 7.3) {
        \archibadge{Client Connection Pool}\archiname{Petasos HA Client Runtime}\\
        \archidesc{Bootstrap: (tcp://pri-a:61616, tcp://bak-a:61617, tcp://pri-b:61618, tcp://bak-b:61619)?ha=true}
    };

    % Replication Group A Nodes
    \node[archi-deployment-node, minimum width=6.5cm] (pri_a) at (3.2, 4.4) {
        \archibadge{Live Primary}\archiname{Artemis Primary A}\\
        \archidesc{Host: 61616, Web: 8161 | Storage: /vol/artemis-pa}
    };
    \node[archi-deployment-node, minimum width=6.5cm] (bak_a) at (3.2, 1.8) {
        \archibadge{Passive Backup}\archiname{Artemis Backup A}\\
        \archidesc{Host: 61617, Web: 8162 | Storage: /vol/artemis-ba}
    };

    % Replication Group B Nodes
    \node[archi-deployment-node, minimum width=6.5cm] (pri_b) at (11.8, 4.4) {
        \archibadge{Live Primary}\archiname{Artemis Primary B}\\
        \archidesc{Host: 61618, Web: 8163 | Storage: /vol/artemis-pb}
    };
    \node[archi-deployment-node, minimum width=6.5cm] (bak_b) at (11.8, 1.8) {
        \archibadge{Passive Backup}\archiname{Artemis Backup B}\\
        \archidesc{Host: 61619, Web: 8164 | Storage: /vol/artemis-bb}
    };

    % --- Connections ---
    % Client connections to primaries
    \draw[archi-serving] (pri_a) -- node[archi-lbl]{Active Core TCP} (pri_a |- client_pool.south);
    \draw[archi-serving] (pri_b) -- node[archi-lbl]{Active Core TCP} (pri_b |- client_pool.south);

    % Replication Channels (Group A & Group B)
    \draw[archi-flow] (pri_a) -- node[archi-lbl]{Replication group-a} (bak_a);
    \draw[archi-flow] (bak_a) -- node[archi-lbl]{Failback Resync} (pri_a);

    \draw[archi-flow] (pri_b) -- node[archi-lbl]{Replication group-b} (bak_b);
    \draw[archi-flow] (bak_b) -- node[archi-lbl]{Failback Resync} (pri_b);

    % Server-Side Cluster Mesh
    \draw[archi-triggering] (pri_a) -- node[archi-lbl]{ON\_DEMAND Cluster Mesh (Redistribution Delay = 0)} (pri_b);
    \draw[archi-triggering] (pri_b) -- (pri_a);

\end{tikzpicture}

    \end{adjustbox}
    \caption{ArchiMate Physical Deployment \& Multi-Broker ActiveMQ Artemis HA Architecture}
    \label{fig:deployment_ha}
\end{figure}

\section{Clustering vs. High Availability (HA) Replication}

Harmonia establishes a clear architectural distinction between horizontal clustering and instance-level replication:

\subsection{1. Server-Side Clustering (\texttt{petasos-cluster})}
\begin{itemize}
    \item \textbf{Active-Active Mesh}: Unites multiple live brokers (\texttt{Primary A} and \texttt{Primary B}) into a distributed message routing fabric.
    \item \textbf{Dynamic Load Distribution (\texttt{ON\_DEMAND})}: Messages published to \texttt{Primary A} are dynamically forwarded across the cluster bridge to \texttt{Primary B} if consumers exist on \texttt{Primary B} and not on \texttt{Primary A}.
    \item \textbf{Immediate Message Redistribution (\texttt{redistribution-delay = 0})}: When all consumers on one broker disconnect, unconsumed messages are immediately redistributed to alternative nodes with active consumers, preventing queue starvation.
    \item \textbf{Horizontal Scalability}: Additional Primary/Backup pairs can be introduced dynamically without disrupting existing message flows.
\end{itemize}

\subsection{2. Journal Replication Groups (\texttt{group-a}, \texttt{group-b})}
\begin{itemize}
    \item \textbf{NIO Journal Streaming}: An active Primary continuously streams file journal updates over TCP to its paired passive Backup replica.
    \item \textbf{Activation on Failure}: If \texttt{Primary A} terminates or loses network connectivity, \texttt{Backup A} detects connection loss within 10 seconds and promotes itself to become the live active broker for \texttt{group-a}.
    \item \textbf{Automatic Failback (\texttt{allow-failback = true})}: When the failed Primary recovers, it connects to the active Backup, resynchronizes any journal deltas written during downtime, and gracefully resumes the Primary role.
\end{itemize}

\section{Isolated Persistent Storage Ownership}

To eliminate single points of failure and storage lock contention, multiple brokers \textbf{never share journal directories or filesystems}. Each broker instance owns an independent persistent Docker volume:

\begin{table}[htbp]
\centering
\small
\begin{tabularx}{\textwidth}{l l p{4.0cm} X}
\toprule
\textbf{Instance Name} & \textbf{HA Role} & \textbf{Docker Volume} & \textbf{Isolated Storage Path} \\
\midrule
\texttt{artemis-primary-a} & Live Primary A & \texttt{artemis\_primary\_a\_data} & \texttt{/var/lib/artemis/data} \\
\texttt{artemis-backup-a} & Replicating Backup A & \texttt{artemis\_backup\_a\_data} & \texttt{/var/lib/artemis/data} \\
\texttt{artemis-primary-b} & Live Primary B & \texttt{artemis\_primary\_b\_data} & \texttt{/var/lib/artemis/data} \\
\texttt{artemis-backup-b} & Replicating Backup B & \texttt{artemis\_backup\_b\_data} & \texttt{/var/lib/artemis/data} \\
\bottomrule
\end{tabularx}
\caption{ActiveMQ Artemis Broker Storage Isolation and Persistent Volume Mapping}
\label{tab:storage_isolation}
\end{table}

\section{Client Reconnection \& Failover Configuration}

All Harmonia modules connect to Petasos using a multi-endpoint high-availability bootstrap URL:
\begin{lstlisting}[language=bash, caption={Petasos Multi-Broker HA Client Bootstrap URL}]
(tcp://artemis-primary-a:61616,tcp://artemis-backup-a:61617,tcp://artemis-primary-b:61618,tcp://artemis-backup-b:61619)?ha=true&reconnectAttempts=-1&retryInterval=500&maxRetryInterval=2000&retryIntervalMultiplier=1.5&connectionTTL=60000&clientFailureCheckPeriod=10000
\end{lstlisting}

\begin{itemize}
    \item \texttt{ha=true}: Enables automatic cluster topology discovery. The client continuously receives live node lists from the brokers.
    \item \texttt{reconnectAttempts=-1}: Directs the client to attempt reconnection indefinitely during transient network partitions.
    \item \texttt{retryInterval=500} \& \texttt{maxRetryInterval=2000}: Exponential backoff with jitter prevents thundering herd reconnection storms on broker reboot.
    \item \texttt{failoverListener}: The \texttt{PetasosArtemisAdapter} registers failover lifecycle listeners to record reconnect telemetry and ensure in-flight transaction integrity.
\end{itemize}

\section{In-Memory Cache Grid HA Topology (\texttt{Mneme})}

The \texttt{Mneme} distributed cache grid deploys across two clustered Infinispan nodes:
\begin{itemize}
    \item \textbf{Replication Strategy}: Caches are configured with \texttt{REPL\_SYNC} mode across \texttt{mneme-node-1} (port \texttt{11222}) and \texttt{mneme-node-2} (port \texttt{11223}).
    \item \textbf{HotRod Client Smart Routing}: Clients maintain TCP connections to both nodes, automatically routing read/write operations to the nearest available replica.
    \item \textbf{Node Crash Resilience}: If either node crashes, the surviving node serves read and write traffic without data loss or interruption.
\end{itemize}

\section{Docker Compose Multi-Container Deployment}

The complete Harmonia platform stack is orchestrated via a unified \texttt{docker-compose.yml} containing:
\begin{itemize}
    \item 4 PostgreSQL 16 relational database containers (2 Clinical nodes, 2 Operations nodes) with persistent isolated volumes.
    \item 4 Spring Boot Mnemosyne JPA persistence servers (2 Clinical HAPI FHIR nodes, 2 Operations nodes).
    \item 2 Infinispan 15+ distributed cache grid nodes (\texttt{infinispan-1}, \texttt{infinispan-2}).
    \item 5 WildFly Jakarta EE 10 application containers (\texttt{pylai-mllp-in}, \texttt{pylai-mllp-out-his}, \texttt{pylai-mllp-out-lis}, \texttt{ponos}, \texttt{iris-befe}).
    \item 2 Vue 3 / Nginx frontend web containers (\texttt{iris-clinical}, \texttt{iris-console}).
\end{itemize}

Health checks with configurable start periods ensure that dependent containers only boot once their prerequisite messaging, caching, and database services report operational readiness.

\subsection{Egress Failure Domain Isolation}
The multi-instance outbound architecture guarantees complete failure domain isolation across external clinical systems. Each egress instance (\texttt{mllp-outbound-his}, \texttt{mllp-outbound-lis}) consumes from its own dedicated Petasos queue (e.g. \texttt{petasos.queue.mllp.outbound.his}). If a destination hospital system experiences downtime, network unreachability, or socket timeouts:
\begin{itemize}
    \item Egress queue backlog is strictly contained within that destination's specific queue without affecting inbound processing or alternate outbound destinations.
    \item Thread pools and socket resources in the faulting gateway instance remain isolated, preventing thread starvation or memory leaks from propagating to other services.
    \item Ponos task sequences can continue dispatching messages to healthy clinical destinations concurrently.
\end{itemize}

\newpage
